% ============================================================================
% CAPÍTULO 5 - GRAFO MULTI-AGENTE COM LANGGRAPH
% ============================================================================

\chapter{Grafo Multi-Agente com LangGraph}\label{cap:multiagente}

Este capítulo apresenta a arquitetura central do assistente: um \textbf{grafo multi-agente} implementado com LangGraph, que orquestra 7 nós especializados para processar consultas médicas de forma segura, rastreável e fundamentada em evidências.

\section{Motivação}

Sistemas de IA para saúde exigem mais do que uma chamada direta a um LLM. É necessário:

\begin{itemize}
    \item \textbf{Classificar} a intenção do usuário antes de agir;
    \item \textbf{Recuperar} dados relevantes (protocolos clínicos e dados do paciente);
    \item \textbf{Validar} a resposta gerada contra regras de segurança;
    \item \textbf{Rastrear} cada etapa para auditoria e explicabilidade.
\end{itemize}

Essas responsabilidades são naturalmente separáveis em agentes especializados, orquestrados por um grafo de estados. O LangGraph, extensão do LangChain, oferece primitivas para construir exatamente esse tipo de sistema: \texttt{StateGraph}, nós que transformam estado, arestas condicionais e compilação para execução.

\section{Visão Geral da Arquitetura}

O grafo possui \textbf{7 nós} (agentes), conectados por \textbf{4 arestas fixas} e \textbf{3 arestas condicionais}. A Figura~\ref{fig:grafo_multiagente} apresenta a topologia completa.

\begin{figure}[H]
\centering
\includegraphics[width=\textwidth]{grafo-multiagente.png}
\caption{Grafo multi-agente LangGraph do Assistente Virtual Médico. Losango representa decisão (router), retângulos representam agentes, elipses representam início/fim. Linhas tracejadas indicam fluxos condicionais.}
\label{fig:grafo_multiagente}
\end{figure}

\subsection{Fluxos Possíveis}

O grafo suporta cinco fluxos distintos, determinados pela classificação da triagem:

\begin{table}[H]
\centering
\caption{Fluxos de execução do grafo multi-agente.}
\label{tab:fluxos_grafo}
\rowcolors{2}{fiap-light}{white}
\small
\begin{tabularx}{\textwidth}{l X}
\toprule
\textbf{Cenário} & \textbf{Caminho no Grafo} \\
\midrule
Protocolo & triagem $\rightarrow$ protocolo $\rightarrow$ raciocínio $\rightarrow$ guardrails $\rightarrow$ explicabilidade $\rightarrow$ logger \\
Paciente & triagem $\rightarrow$ paciente\_data $\rightarrow$ raciocínio $\rightarrow$ guardrails $\rightarrow$ explicabilidade $\rightarrow$ logger \\
Ambos & triagem $\rightarrow$ paciente\_data $\rightarrow$ protocolo $\rightarrow$ raciocínio $\rightarrow$ guardrails $\rightarrow$ explicabilidade $\rightarrow$ logger \\
Fora de escopo & triagem $\rightarrow$ logger $\rightarrow$ END \\
Retry & $\ldots$ $\rightarrow$ guardrails $\rightarrow$ raciocínio $\rightarrow$ guardrails $\rightarrow$ $\ldots$ (máx.\ 2$\times$) \\
\bottomrule
\end{tabularx}
\end{table}

\section{State Schema}\label{sec:state}

O estado compartilhado entre todos os nós é definido pelo \texttt{MedicalAssistantState}, um \texttt{TypedDict} com 15 variáveis organizadas em cinco grupos:

\begin{table}[H]
\centering
\caption{Variáveis do \texttt{MedicalAssistantState}.}
\label{tab:state_schema}
\rowcolors{2}{fiap-light}{white}
\small
\begin{tabularx}{\textwidth}{l l l X}
\toprule
\textbf{Grupo} & \textbf{Variável} & \textbf{Tipo} & \textbf{Descrição} \\
\midrule
\multirow{2}{*}{Entrada}
    & \texttt{query} & \texttt{str} & Pergunta do médico \\
    & \texttt{patient\_id} & \texttt{Optional[str]} & ID ou nome do paciente \\
\midrule
\multirow{2}{*}{Triagem}
    & \texttt{query\_type} & \texttt{Literal[4]} & protocolo $|$ paciente $|$ ambos $|$ fora\_de\_escopo \\
    & \texttt{entities} & \texttt{dict} & Entidades extraídas (paciente, condição, etc.) \\
\midrule
\multirow{2}{*}{Contexto}
    & \texttt{protocols} & \texttt{list[dict]} & Chunks recuperados do ChromaDB \\
    & \texttt{patient\_data} & \texttt{Optional[dict]} & Dados clínicos do PostgreSQL \\
\midrule
\multirow{4}{*}{Raciocínio}
    & \texttt{draft\_response} & \texttt{str} & Rascunho gerado pelo LLM \\
    & \texttt{guardrail\_result} & \texttt{Literal[2]} & aprovado $|$ reprovado \\
    & \texttt{guardrail\_feedback} & \texttt{Optional[str]} & Motivo da reprovação \\
    & \texttt{retry\_count} & \texttt{int} & Contador de retries (máx.\ 2) \\
\midrule
\multirow{4}{*}{Saída}
    & \texttt{final\_response} & \texttt{str} & Resposta formatada com citações \\
    & \texttt{sources} & \texttt{list[str]} & Lista de fontes citadas \\
    & \texttt{confidence} & \texttt{Literal[3]} & alta $|$ media $|$ baixa \\
    & \texttt{warnings} & \texttt{list[str]} & Advertências ao médico \\
\midrule
Auditoria
    & \texttt{audit\_log} & \texttt{list[dict]} & Trace completo da execução \\
\bottomrule
\end{tabularx}
\end{table}

\begin{lstlisting}[caption={Definição do \texttt{MedicalAssistantState} em \texttt{src/flows/state.py}.}]
class MedicalAssistantState(TypedDict):
    # Entrada
    query: str
    patient_id: Optional[str]
    # Triagem
    query_type: Literal["protocolo", "paciente",
                        "ambos", "fora_de_escopo"]
    entities: dict
    # Contexto
    protocols: list[dict]
    patient_data: Optional[dict]
    # Raciocinio
    draft_response: str
    guardrail_result: Literal["aprovado", "reprovado"]
    guardrail_feedback: Optional[str]
    retry_count: int
    # Saida
    final_response: str
    sources: list[str]
    confidence: Literal["alta", "media", "baixa"]
    warnings: list[str]
    # Auditoria
    audit_log: list[dict]
\end{lstlisting}

\section{Os 7 Nós do Grafo}\label{sec:agentes}

Cada nó é uma função Python que recebe o estado atual, executa sua lógica e retorna um dicionário parcial com as variáveis que deseja atualizar. O LangGraph faz o merge automático.

\subsection{Triagem (Router)}

O agente de triagem é o \textbf{ponto de entrada} do grafo. Recebe a query do médico e:

\begin{enumerate}
    \item \textbf{Classifica} o tipo da consulta em uma de quatro categorias;
    \item \textbf{Extrai entidades} relevantes (nome do paciente, condição médica, medicamento, exame).
\end{enumerate}

Quando a LLM está disponível, a classificação é feita via prompt estruturado que retorna JSON. Caso contrário, um classificador heurístico baseado em palavras-chave é utilizado como fallback.

\begin{table}[H]
\centering
\caption{Classificação da triagem e roteamento.}
\label{tab:triagem_routing}
\rowcolors{2}{fiap-light}{white}
\begin{tabularx}{\textwidth}{l X l}
\toprule
\textbf{Tipo} & \textbf{Critério} & \textbf{Destino} \\
\midrule
\texttt{protocolo} & Pergunta sobre tratamento, diagnóstico ou conduta clínica genérica & Protocolo \\
\texttt{paciente} & Pergunta sobre dados de um paciente específico & Paciente Data \\
\texttt{ambos} & Protocolo clínico aplicado a um paciente específico & Paciente Data \\
\texttt{fora\_de\_escopo} & Pergunta não relacionada à medicina & Logger \\
\bottomrule
\end{tabularx}
\end{table}

\begin{infobox}[Ajuste Automático de Classificação]
Quando o campo \texttt{patient\_id} é fornecido na entrada, dois ajustes automáticos são aplicados:
\begin{itemize}
    \item Se a triagem classifica como \texttt{protocolo} $\rightarrow$ reclassifica para \texttt{ambos};
    \item Se classifica como \texttt{paciente} mas a query contém \textbf{sintomas clínicos} (dor, febre, palpitações, dispneia, etc.) $\rightarrow$ reclassifica para \texttt{ambos}, garantindo busca de protocolos relevantes aos sintomas.
\end{itemize}
Além disso, o classificador heurístico detecta palavras-chave de sintomas (``dor no peito'', ``palpitações'', ``febre'') e classifica diretamente como \texttt{ambos} quando coexistem com referência a paciente.
\end{infobox}

\subsection{Protocolo (RAG ChromaDB)}

O agente de protocolo realiza \textbf{busca semântica} nos 544 documentos indexados no ChromaDB (14 Linhas de Cuidado + 5 protocolos de emergência, conforme Capítulo~\ref{cap:rag}).

\begin{itemize}
    \item \textbf{Limpeza da query:} Antes da busca semântica, a query é pré-processada: o nome do paciente é removido e \textit{stop words} em português são filtradas, mantendo apenas os termos clínicos relevantes (ex: ``A paciente Maria Alice está com dor no peito'' $\rightarrow$ ``dor peito'');
    \item \textbf{Busca híbrida (sintomas + comorbidades):} Quando dados do paciente estão disponíveis (fluxo ``ambos''), o agente executa três camadas de busca:
    \begin{enumerate}
        \item \textbf{Query base} (termos clínicos limpos) --- captura protocolos de emergência relevantes aos sintomas mencionados;
        \item \textbf{Uma busca por comorbidade} (ex: ``dor peito + Hipertensão'', ``dor peito + Diabetes tipo 2'') --- garante cobertura equilibrada das condições crônicas do paciente.
    \end{enumerate}
    Sem dados do paciente, a busca é ampla com a query limpa;
    \item \textbf{Ranking por relevância:} Após a recuperação, os resultados são classificados como \texttt{direta} (contêm termos das comorbidades) ou \texttt{complementar} (informação geral), ordenados por relevância;
    \item Retorna conteúdo, fonte, seção, distância e relevância para cada resultado.
\end{itemize}

\begin{infobox}[Por que busca híbrida?]
A busca por comorbidade garante cobertura de todas as condições crônicas do paciente, mas não captura protocolos de emergência associados aos \textbf{sintomas} da consulta. Por exemplo, ``dor no peito e palpitações'' deve encontrar o protocolo de IAM (infarto), mesmo que o paciente não tenha essa comorbidade. A busca híbrida --- primeira camada pelos sintomas, segunda pelas comorbidades --- resolve esse problema. A limpeza de stop words é necessária porque frases longas (``A paciente está se queixando de\ldots'') diluem a busca semântica, favorecendo resultados irrelevantes.
\end{infobox}

\subsection{Paciente Data (PostgreSQL)}

O agente de dados do paciente consulta o banco PostgreSQL via SQLAlchemy, buscando:

\begin{itemize}
    \item \textbf{Dados cadastrais:} nome, sexo, data de nascimento, comorbidades, alergias, medicamentos em uso;
    \item \textbf{Exames:} tipo, data, resultados e interpretação (últimos 5);
    \item \textbf{Prontuários:} data, diagnóstico CID, texto diagnóstico (últimos 3);
    \item \textbf{Receitas:} data, condição, medicamentos prescritos.
\end{itemize}

A busca é feita por nome (parcial, case-insensitive) ou por ID do paciente. Todos os valores são serializados para JSON, tratando tipos como \texttt{date} e \texttt{UUID}.

\subsection{Raciocínio (Mistral 7B Fine-Tuned)}

O agente de raciocínio é o \textbf{cérebro} do sistema. Utiliza o modelo Mistral 7B com fine-tuning (Capítulo~3) para sintetizar uma resposta clínica a partir de três fontes de informação:

\begin{enumerate}
    \item A \textbf{query} original do médico;
    \item Os \textbf{protocolos clínicos} recuperados pelo agente de protocolo;
    \item Os \textbf{dados do paciente} recuperados pelo agente de dados.
\end{enumerate}

O prompt é montado dinamicamente com as seções de contexto disponíveis. Em caso de \textbf{retry} (quando o guardrails reprova a resposta), o feedback de reprovação é incluído no prompt para que o modelo corrija os problemas apontados.

\begin{lstlisting}[caption={Estrutura simplificada do prompt do agente de raciocínio.}]
Voce e um assistente medico especializado. Com base no
contexto fornecido, responda a pergunta do medico.

PERGUNTA DO MEDICO: {query}

PROTOCOLOS CLINICOS RELEVANTES:
[1] Fonte: Pathway-Hipertensao | Secao: Tratamento
    Os principios que norteiam o tratamento...

DADOS DO PACIENTE:
Nome: Maria Santos | Sexo: F | Nascimento: 1958-03-15
Comorbidades: hipertensao, diabetes tipo 2
Ultimos exames:
  - HbA1C (2026-01-15): 7.2% - controle moderado

# Se retry:
ATENCAO - SUA RESPOSTA ANTERIOR FOI REPROVADA:
Motivo: Prescricao direta detectada. Reformule como
sugestao e indique validacao medica.

RESPOSTA:
\end{lstlisting}

Quando o modelo fine-tuned não está disponível, um sistema de \textbf{fallback} gera respostas baseadas em template com o contexto recuperado. O fallback apresenta os dados completos do paciente (comorbidades, exames, prontuários, receitas) seguidos de até \textbf{6 protocolos distintos}, deduplicados por fonte e seção, com conteúdo truncado no final da última sentença completa. Protocolos complementares são indicados com label, e o limite de 6 garante cobertura de todas as fontes relevantes sem repetição.

\subsection{Guardrails (Segurança)}

O agente de guardrails aplica \textbf{4 regras de segurança} obrigatórias a cada resposta:

\begin{table}[H]
\centering
\caption{Regras de validação dos guardrails.}
\label{tab:guardrails_rules}
\rowcolors{2}{fiap-light}{white}
\begin{tabularx}{\textwidth}{l l X}
\toprule
\textbf{Regra} & \textbf{Tipo} & \textbf{O que verifica} \\
\midrule
\texttt{no\_prescription} & Regex & Detecta padrões de prescrição direta (``prescrevo'', ``tome 500mg'', ``administre'') \\
\texttt{human\_validation} & Keyword & Verifica presença de disclaimer de validação humana \\
\texttt{no\_definitive\_diagnosis} & Regex & Detecta diagnósticos definitivos (``o paciente tem'', ``confirmo que'') \\
\texttt{context\_consistency} & Heurística & Verifica se resposta não é excessivamente curta dado o contexto \\
\bottomrule
\end{tabularx}
\end{table}

\subsubsection{Mecanismo de Retry}

Quando uma regra falha, o guardrails retorna \texttt{reprovado} com feedback detalhado. O grafo então reenvia o estado ao agente de raciocínio, que recebe o feedback no prompt e gera uma resposta corrigida. Este ciclo pode ocorrer no \textbf{máximo 2 vezes}. Caso o limite seja atingido, a resposta segue com advertências.

\begin{figure}[H]
\centering
\begin{tikzpicture}[
    node distance=1cm and 1.5cm,
    block/.style={rectangle, draw, rounded corners=3pt, minimum width=2.8cm, minimum height=0.9cm, align=center, font=\small\sffamily},
    decision/.style={diamond, draw, minimum width=2cm, minimum height=1.5cm, align=center, font=\small\sffamily, inner sep=1pt},
    process/.style={block, fill=fiap-red!15, draw=fiap-red, thick},
    safe/.style={block, fill=green!10, draw=green!60!black, thick},
    warn/.style={block, fill=orange!10, draw=orange!70, thick},
    arrow/.style={->, thick, >=stealth}
]

\node[process] (raciocinio) {Raciocínio\\(gera draft)};
\node[warn, right=2cm of raciocinio] (guardrails) {Guardrails\\(4 regras)};
\node[decision, right=2cm of guardrails] (decision) {Aprovado?};
\node[safe, above right=0.8cm and 1.5cm of decision] (explicabilidade) {Explicabilidade};
\node[process, below right=0.8cm and 1.5cm of decision] (retry) {Retry\\(com feedback)};

\draw[arrow] (raciocinio) -- (guardrails);
\draw[arrow] (guardrails) -- (decision);
\draw[arrow] (decision) -- node[above, font=\tiny\sffamily] {Sim} (explicabilidade);
\draw[arrow] (decision) -- node[below, font=\tiny\sffamily] {Não (máx 2x)} (retry);
\draw[arrow] (retry.south) -- ++(0,-0.5) -| node[below, pos=0.25, font=\tiny\sffamily] {feedback no prompt} (raciocinio.south);

\end{tikzpicture}
\caption{Mecanismo de retry dos guardrails: respostas reprovadas são devolvidas ao raciocínio com feedback.}
\label{fig:guardrails_retry}
\end{figure}

\subsection{Explicabilidade (Citações e Confiança)}

O agente de explicabilidade é responsável por tornar a resposta \textbf{transparente e rastreável}:

\begin{enumerate}
    \item \textbf{Extrai citações} dos protocolos e dados utilizados:
    \begin{itemize}
        \item Formato: \texttt{[Pathway-Diabetes, Seção: Tratamento]}
        \item Formato: \texttt{[Exame: HbA1C, Data: 2026-01-15]}
    \end{itemize}

    \item \textbf{Avalia confiança} com base em múltiplos fatores:

    \begin{table}[H]
    \centering
    \caption{Critérios de avaliação de confiança.}
    \label{tab:confianca}
    \rowcolors{2}{fiap-light}{white}
    \begin{tabular}{l c}
    \toprule
    \textbf{Fator} & \textbf{Pontuação} \\
    \midrule
    3+ protocolos recuperados & +3 \\
    1-2 protocolos recuperados & +2 \\
    Distância média $<$ 0.5 & +2 \\
    Distância média $<$ 1.0 & +1 \\
    Dados do paciente disponíveis & +2 \\
    Cada retry do guardrails & $-1$ \\
    \midrule
    $\geq 5$ pontos: \textbf{alta} & $\geq 3$: \textbf{média} \\
    \multicolumn{2}{c}{$< 3$ pontos: \textbf{baixa}} \\
    \bottomrule
    \end{tabular}
    \end{table}

    \item \textbf{Gera advertências} quando apropriado (confiança baixa, sem protocolos, retries realizados);

    \item \textbf{Separa fontes por relevância:} fontes \textbf{diretas} (relacionadas ao quadro do paciente) e \textbf{complementares} (informação geral) são apresentadas em seções distintas;

    \item \textbf{Formata a resposta final} com seções: resposta, fontes diretas, fontes complementares, confiança, advertências e disclaimer obrigatório.
\end{enumerate}

\subsection{Logger (Auditoria)}

O agente de logger é o \textbf{nó final} em todos os fluxos. Consolida o estado completo em um registro JSON estruturado e persiste em arquivo diário (\texttt{logs/audit\_YYYY-MM-DD.jsonl}).

\begin{lstlisting}[caption={Estrutura do registro de auditoria.}]
{
    "timestamp": "2026-03-06T18:30:00",
    "session": {
        "query": "Qual tratamento para hipertensao?",
        "patient_id": null
    },
    "triagem": {
        "query_type": "protocolo",
        "entities": {"condicao": "hipertensao"}
    },
    "context": {
        "protocols_count": 5,
        "protocols_sources": ["Pathway-Hipertensao", ...],
        "patient_data_available": false
    },
    "response": {
        "confidence": "alta",
        "sources": ["[Pathway-Hipertensao, ...]"],
        "warnings": []
    },
    "guardrails": {
        "result": "aprovado",
        "retry_count": 0
    },
    "agent_trace": [
        {"agent": "triagem", "timestamp": "..."},
        {"agent": "protocolo", "timestamp": "..."},
        ...
    ]
}
\end{lstlisting}

Para consultas classificadas como \texttt{fora\_de\_escopo}, o logger gera uma resposta padrão educada informando o escopo do assistente, sem invocar nenhum outro agente.

\section{Montagem do Grafo}\label{sec:montagem}

A montagem do grafo é feita em \texttt{src/flows/graph.py} utilizando a API do LangGraph:

\begin{lstlisting}[caption={Montagem do grafo em \texttt{src/flows/graph.py}.}]
from langgraph.graph import StateGraph, END

graph = StateGraph(MedicalAssistantState)

# 7 nos
graph.add_node("triagem", triagem_agent)
graph.add_node("protocolo", protocolo_agent)
graph.add_node("paciente_data", paciente_data_agent)
graph.add_node("raciocinio", raciocinio_agent)
graph.add_node("guardrails", guardrails_agent)
graph.add_node("explicabilidade", explicabilidade_agent)
graph.add_node("logger", logger_agent)

# Entry point
graph.set_entry_point("triagem")

# Arestas condicionais
graph.add_conditional_edges("triagem",
    route_after_triagem, {
    "protocolo": "protocolo",
    "paciente_data": "paciente_data",
    "logger": "logger",
})
graph.add_conditional_edges("paciente_data",
    route_after_paciente_data, {
    "protocolo": "protocolo",
    "raciocinio": "raciocinio",
})
graph.add_conditional_edges("guardrails",
    route_after_guardrails, {
    "explicabilidade": "explicabilidade",
    "raciocinio": "raciocinio",
})

# Arestas fixas
graph.add_edge("protocolo", "raciocinio")
graph.add_edge("raciocinio", "guardrails")
graph.add_edge("explicabilidade", "logger")
graph.add_edge("logger", END)

app = graph.compile()
\end{lstlisting}

\section{Funções de Roteamento}

As três funções de roteamento condicional determinam o caminho de execução no grafo com base no estado atual:

\begin{table}[H]
\centering
\caption{Funções de roteamento condicional.}
\label{tab:routing_functions}
\rowcolors{2}{fiap-light}{white}
\small
\begin{tabularx}{\textwidth}{l l X}
\toprule
\textbf{Função} & \textbf{Após} & \textbf{Lógica} \\
\midrule
\texttt{route\_after\_triagem} & Triagem & Se \texttt{protocolo} $\rightarrow$ protocolo; se \texttt{paciente} ou \texttt{ambos} $\rightarrow$ paciente\_data; senão $\rightarrow$ logger \\
\texttt{route\_after\_paciente\_data} & Paciente Data & Se \texttt{ambos} $\rightarrow$ protocolo (busca contextualizada); senão $\rightarrow$ raciocínio \\
\texttt{route\_after\_guardrails} & Guardrails & Se \texttt{aprovado} $\rightarrow$ explicabilidade; se \texttt{reprovado} e retry $<$ 2 $\rightarrow$ raciocínio; senão $\rightarrow$ explicabilidade (com warnings) \\
\bottomrule
\end{tabularx}
\end{table}

\section{Exemplo de Execução End-to-End}

Para ilustrar o funcionamento completo, considere a seguinte consulta:

\begin{warningbox}[Consulta de Exemplo]
\textbf{Médico:} ``Quais exames de seguimento pedir para a paciente Maria Santos, diabética tipo 2?''
\end{warningbox}

\begin{enumerate}
    \item \textbf{Triagem} $\rightarrow$ \texttt{query\_type="ambos"}, \texttt{entities=\{paciente: "Maria Santos", condicao: "diabetes"\}}
    \item \textbf{Paciente Data} $\rightarrow$ PostgreSQL $\rightarrow$ Maria Santos: comorbidades=[hipertensão, diabetes], exames, medicamentos
    \item \textbf{Protocolo} $\rightarrow$ Busca híbrida: (1) query limpa ``exames seguimento diabética'' $\rightarrow$ protocolos de diabetes e emergência; (2) ``+ hipertensão'' e ``+ diabetes'' por comorbidade $\rightarrow$ Pathway-Hipertensão + Pathway-Diabetes (diretos) + protocolos complementares
    \item \textbf{Raciocínio} $\rightarrow$ Mistral 7B FT sintetiza resposta com protocolos relevantes + dados da paciente
    \item \textbf{Guardrails} $\rightarrow$ 4 regras verificadas $\rightarrow$ \texttt{aprovado}
    \item \textbf{Explicabilidade} $\rightarrow$ Confiança: alta | Fontes diretas: [Pathway-Diabetes, Pathway-Hipertensão, Exame HbA1C]
    \item \textbf{Logger} $\rightarrow$ Registro JSON persistido em \texttt{logs/audit\_2026-03-06.jsonl}
\end{enumerate}

\section{Estrutura de Arquivos}

\begin{table}[H]
\centering
\caption{Arquivos do módulo \texttt{src/flows/}.}
\label{tab:arquivos_flows}
\rowcolors{2}{fiap-light}{white}
\small
\begin{tabularx}{\textwidth}{l X}
\toprule
\textbf{Arquivo} & \textbf{Responsabilidade} \\
\midrule
\texttt{state.py} & Definição do \texttt{MedicalAssistantState} \\
\texttt{graph.py} & Montagem do grafo, funções de roteamento e \texttt{run\_assistant()} \\
\texttt{agents/triagem.py} & Classificação da query e extração de entidades \\
\texttt{agents/protocolo.py} & Busca semântica no ChromaDB \\
\texttt{agents/paciente\_data.py} & Consulta SQL no PostgreSQL \\
\texttt{agents/raciocinio.py} & Síntese de resposta via LLM fine-tuned \\
\texttt{agents/guardrails.py} & 4 regras de validação de segurança \\
\texttt{agents/explicabilidade.py} & Citações, confiança e formatação final \\
\texttt{agents/logger\_agent.py} & Consolidação e persistência de auditoria \\
\bottomrule
\end{tabularx}
\end{table}

\section{Decisões de Design}

\begin{enumerate}
    \item \textbf{Estado compartilhado vs.\ mensagens:} Optamos por estado compartilhado (\texttt{TypedDict}) ao invés de troca de mensagens entre agentes. Isso simplifica o fluxo e permite que cada agente acesse o contexto completo sem repropagação.

    \item \textbf{Fallbacks em cada agente:} Cada agente possui lógica de fallback para funcionar mesmo sem dependências externas (LLM, ChromaDB, PostgreSQL), facilitando testes unitários e desenvolvimento incremental.

    \item \textbf{Retry com feedback:} Ao invés de simplesmente rejeitar respostas inseguras, o mecanismo de retry reinjeta o feedback do guardrails no prompt, dando à LLM a oportunidade de se autocorrigir.

    \item \textbf{Auditoria acumulativa:} O campo \texttt{audit\_log} é incrementado por cada agente (append-only), criando um trace completo de execução sem overhead de um sistema de observabilidade externo.

    \item \textbf{LLM injetável:} A LLM é passada via parâmetro (\texttt{run\_assistant(query, llm=...)}) ao invés de ser instanciada dentro do grafo, permitindo trocar facilmente entre modelo base, fine-tuned ou mock para testes. O carregador (\texttt{\_load\_llm()} em \texttt{app/main.py}) aceita tanto um path local (\texttt{models/medical-assistant-final}) quanto um repositório do Hugging Face Hub (\texttt{zagari/medical-assistant-mistral-7b-ft}), configurável pela variável de ambiente \texttt{FINE\_TUNED\_MODEL}. Quando o modelo não está disponível, o agente de raciocínio utiliza o fallback baseado em template.

    \item \textbf{Paciente antes de protocolo + busca híbrida:} Para queries tipo ``ambos'', o grafo busca os dados do paciente \emph{antes} dos protocolos. Isso permite busca híbrida no ChromaDB: primeiro pelos \textbf{sintomas} da query (capturando protocolos de emergência como IAM quando há ``dor no peito''), depois por \textbf{comorbidade} do paciente (cobertura de condições crônicas). A query é pré-processada com remoção de stop words e nome do paciente para maximizar a relevância semântica.
\end{enumerate}
